% Standalone formula guide, companion volume: Chapter 5.
% Source: contents/chapter5_shortened.tex.
% Compile this file as the main document. Bibliography is embedded.
%
% NO-REPETITION RULE: Chapter 5 contains NO displayed mathematics. Every
% general formula it applies is defined in formula_guide_chapters1_4.tex and is
% referred to here by name and number, never reprinted. The equations below are
% either (a) quantities Chapter 5 states in prose and this guide promotes, or
% (b) numeric evaluations of an earlier formula at specific counts.
\documentclass[11pt,a4paper]{article}
\usepackage[T1]{fontenc}
\usepackage[utf8]{inputenc}
\usepackage[margin=25mm]{geometry}
\usepackage{amsmath,amssymb}
\usepackage[hidelinks]{hyperref}
\usepackage{url}
\numberwithin{equation}{section}
\setlength{\parindent}{0pt}
\setlength{\parskip}{0.5em}
\allowdisplaybreaks
\title{Formula Guide: Chapter 5}
\author{}
\date{}
\begin{document}
\maketitle

\textbf{Scope.} This companion volume covers the shortened results chapter,
\texttt{contents/chapter5\_shortened.tex}. Section and equation numbers
continue those of \texttt{formula\_guide\_chapters1\_4.tex}, so a reference of
the form ``Equation (2.14)'' points into that document.

\textbf{Finding, stated plainly.} \textbf{Chapter 5 contains no displayed
equation blocks and defines no new formula.} This was verified by scanning the
source for displayed-math delimiters and equation environments; the count is
zero for both. What the chapter contains is fourteen inline mathematical
spans, and all but two of those are notation rather than formulas: tensor
shapes, configuration-pair arrows and a difference symbol.

That is what a results chapter should look like. Its work is to report
measured numbers, and the formulas that produce them were defined in
Chapters~2 to~4.

\textbf{What this volume therefore does.} Three things, none of which involves
restating an earlier formula:

\begin{enumerate}
	\item Promotes the two aggregation rules the chapter uses throughout but
	only ever states in prose --- the matched-pair difference and the mean
	marginal effect.
	\item Works through the one genuine calculation the chapter performs, the
	all-Negative reference classifier, as a numeric evaluation of metric
	formulas defined earlier.
	\item Indexes the chapter's notation, and records which earlier equation
	each restated quantity belongs to.
\end{enumerate}

\tableofcontents
\clearpage

\setcounter{section}{4}
\section{Chapter 5: Results}
\label{fg5:chapter5}

\subsection{Aggregation rules used throughout the chapter}

\subsubsection{The matched-pair difference}
\label{fg5:pair-difference}
\textbf{Thesis locations:} 5.2--5.5, every ``The six pairs'' and ``The four
pairs'' table.

The chapter's central quantity is written in its tables as
``$\Delta$ pooled macro F1'' and in its prose as a signed change, but never as
an equation. Promoted, it is
\begin{equation}
    \Delta_j = S\!\left(\mathrm{on}_j\right) - S\!\left(\mathrm{off}_j\right),
    \label{fg5:eq:pair-difference}
\end{equation}
where $j$ indexes a matched pair, $S(\cdot)$ is pooled macro F1 as constructed
in Equation~(2.26) and pooled per Equation~(2.28), and
$\mathrm{on}_j$ and $\mathrm{off}_j$ are the two configurations of pair $j$ ---
identical in every flag except the one under test.

The construction is what gives the chapter its evidential standing. Because
the two members differ in exactly one flag, and because every non-varied
factor, seed and fold assignment is held identical, $\Delta_j$ isolates that
flag within that architectural context. The chapter's arrow notation,
$\texttt{config\_1} \rightarrow \texttt{config\_2}$, names the pair by its off
and on members in that order.

Two cautions the chapter states explicitly. Differences are calculated from
\emph{unrounded} results, so subtracting the displayed four-decimal scores can
disagree in the last digit. And a single seed means $\Delta_j$ is a point
estimate carrying no uncertainty interval; it cannot establish statistical
significance.

\textbf{Source.} Project design. The matched-pair form is the study's own
construction, not an aggregation rule imported from a reviewed paper.

\subsubsection{The mean marginal effect}
\label{fg5:mean-effect}
\textbf{Thesis locations:} 5.2.1, 5.3.1, 5.4.1, 5.5.1, each ``The effect''.

Each component's headline number is the mean of its pair differences:
\begin{equation}
    \overline{\Delta} = \frac1J \sum_{j=1}^{J} \Delta_j,
    \label{fg5:eq:mean-effect}
\end{equation}
with $J = 6$ for the transformer and for EVM, and $J = 4$ for SimAM and for
the convolutional stem. The counts differ because the validity rule of
Equation~(4.1) removes the four cells that pair SimAM with no CNN, so those
two components have fewer architecturally valid pairs to toggle within.

The chapter is careful about what this average can carry, and the guide
repeats that care rather than the algebra. A mean over four or six
single-seed differences has no distributional interpretation: it does not
estimate a population effect, and its spread across pairs is evidence about
architectural dependence rather than sampling noise. Chapter~5's own
subsection ``Why the mean is not the finding'' makes this argument for SimAM,
where the sign varies across pairs.

For the transformer, $\overline{\Delta} = +0.2173$ with all six differences
positive, ranging $+0.1578$ to $+0.2786$. Consistency of sign across six of
six pairs is a stronger observation than the mean itself, because it does not
depend on the magnitudes.

\textbf{Source.} Project design, as above.

\subsection{The all-Negative reference classifier: thesis section 5.1.3}

\subsubsection{Numeric evaluation of the metric definitions}
\label{fg5:all-negative}
\textbf{Thesis location:} 5.1.3, \texttt{sec:all-negative-reference}.

\textbf{No new formula.} This is the chapter's one genuine calculation, and it
evaluates formulas already defined: accuracy via
Equation~(2.27), per-class F1 via Equation~(2.25) and the macro average via
Equation~(2.26). Only the arithmetic is new.

Consider a classifier that ignores its input and predicts Negative for every
one of the 156 clips, 99 of which are Negative. Then
\begin{equation}
\begin{aligned}
    \text{accuracy} &= \frac{99}{156} = 0.6346,\\
    F1_{\mathrm{Negative}} &= \frac{2\times 99}{156+99} = 0.7765,\\
    F1_{\mathrm{macro}} &= \frac{0.7765+0+0}{3} = 0.2588.
    \end{aligned}
    \label{fg5:eq:all-negative}
\end{equation}

The middle line is Equation~(2.25) with $TP = 99$, $FP = 57$ and $FN = 0$:
every Negative clip is found, so there are no false negatives, and every
non-Negative clip is wrongly called Negative, giving
$2TP + FP + FN = 198 + 57 = 255 = 156 + 99$. The third line is the macro
average with the Positive and Surprise scores at zero, since the classifier
never predicts either, and the project convention of Equation~(2.25) assigns
zero rather than omitting an undefined class.

\textbf{What it establishes, and what it does not.} All twelve configurations
exceed 0.2588 on pooled macro F1. That is a real but weak reference: it
rules out one specific degenerate rule, not every input-independent rule, and
it does not establish generalisation. The accuracy line carries the sharper
warning --- \textbf{0.6346 accuracy is attainable without discriminating among
expressions at all}, which is why accuracy is not the primary metric and why
configuration~8's accuracy lead over configuration~2 (0.7500 against 0.7436,
117 against 116 correct clips) does not settle the ranking.

\textbf{Source.} This thesis's own arithmetic over its class distribution. The
metric definitions are attributed at their own equations; See et
al.~\cite{see2019} supply the pooled construction.

\subsubsection{The mean-of-folds ceiling, restated}
\label{fg5:ceiling-ch5}
\textbf{Thesis location:} 5.1.3.

\textbf{No new formula.} Chapter~5 reports a \textbf{0.6267} ceiling on
mean-of-folds macro F1 imposed by fold composition. This is
Equation~(4.4), which Chapter~4 states as $\approx 0.627$; the two agree, the
results chapter simply carrying one more digit. The ceiling applies with all
three classes retained and undefined class F1 set to zero, and it does not
apply to pooled macro F1. Figure 5.2 of the thesis plots the two aggregations
against each other with that ceiling marked.

\textbf{Source.} As Equation~(4.4): this thesis's own arithmetic over its fold
composition.

\subsection{Notation index}
\label{fg5:notation}

The chapter's remaining mathematical spans are notation, not formulas. They
are listed here so the index is complete.

\textbf{$\rightarrow$ (configuration pairs).} In
$\texttt{config\_1} \rightarrow \texttt{config\_2}$ and the abbreviated forms
$2\rightarrow12$, $9\rightarrow7$, $5\rightarrow16$, $13\rightarrow16$,
$9\rightarrow6$ and $2\rightarrow9$, the arrow separates the off member from
the on member of a matched pair, in the sense of
Equation~\eqref{fg5:eq:pair-difference}. The numerals are configuration
identifiers, not ranks or magnitudes.

\textbf{$\Delta$.} The difference of
Equation~\eqref{fg5:eq:pair-difference}, used as a column heading.

\textbf{$-$.} A typeset minus sign on negative differences, written
explicitly so that a leading minus is not mistaken for a hyphen.

\textbf{$(1,3,3)$.} The convolutional kernel of the stem, in the
$(k_T,k_H,k_W)$ convention set out in Chapter 2's entry on
three-dimensional tensor and kernel notation. Because $k_T = 1$, the stem performs \textbf{no temporal
convolution}; Chapter~5 relies on this when it describes the transformer-free
control as order-blind. The control is nevertheless not independent across
frames at every operation, since batch normalisation during training and
SimAM whenever enabled both use statistics spanning time.

\textbf{$224\times224$ and $28\times28$.} Spatial input dimensions in pixels,
in the width-by-height pixel convention of Chapter 2's image-dimensions
entry. The comparison is raised as one of three unresolved candidate
explanations for the stem's weak contribution: this pipeline's
$224\times224$ input is large relative to STSTNet's
$28\times28$~\cite{liong2019b}. The chapter declines to conclude from it,
noting that other reviewed systems also use $224\times224$ and that shallow
models need not share deeper models' resolution sensitivity. No universal
literature threshold is established.

\subsection{Coverage boundary}

All fourteen inline mathematical spans in the shortened Chapter 5 are
accounted for above: two promoted aggregation rules, one numeric evaluation
across three lines, one restated ceiling, and ten notation items. The chapter
has \textbf{no displayed equation blocks}, so none has been omitted, and no
formula from Chapters 1 to 4 has been reprinted here.

Quantities the chapter reports without expressing as a calculation --- pooled
macro F1 scores, peak VRAM, extrapolated GPU-hours, per-class F1 tables and
confusion matrices --- are measurements, not formulas, and have not been
reverse-engineered into equations.

\begin{thebibliography}{9}
\bibitem{see2019}
John See, Moi Hoon Yap, Jingting Li, Xiaopeng Hong and Su-Jing Wang,
``MEGC 2019 --- The Second Facial Micro-Expression Grand Challenge,''
\emph{14th IEEE International Conference on Automatic Face and Gesture
Recognition (FG 2019)}, 2019.

\bibitem{liong2019b}
Sze-Teng Liong, Y. S. Gan, John See, Huai-Qian Khor and Yen-Chang Huang,
``Shallow Triple Stream Three-Dimensional CNN (STSTNet) for Micro-Expression
Recognition,'' \emph{14th IEEE International Conference on Automatic Face and
Gesture Recognition (FG 2019)}, 2019.
\end{thebibliography}

\end{document}
